\end{Theorem}

\begin{proof}

For each tableau $T(R)$, we have an explicit construction of the module containing $T(R)$ (recall
Def\/inition~\ref{def-cont}):
\begin{gather*}
M(T(R)):=U\cdot T(R)/\left(\sum U\cdot T(Q)\right),
\end{gather*}
where the sum is taken over tableaux $T(Q)$ such that $T(Q)\in U\cdot T(R)$ and $U\cdot T(Q)$ is a~proper submodule of
$U\cdot T(R)$.

The module $M(T(R))$ is simple.
Indeed, this follows from the fact that for any nonzero tableau $T(S)$ in $M(T(R))$ we have $U\cdot T(S)=U\cdot T(R)$
and, hence, $T(S)$ generates $M(T(R))$.

By Theorem~\ref{basis of submodule generated by T} and Proposition~\ref{when two tableaux generate the same submodule},
a~basis for a~proper submodule $U\cdot T(Q)$ of $U\cdot T(R)$ is
$\{T(S):\Omega^{+}(T(R))\subsetneq\Omega^{+}(T(Q))\subseteq \Omega^{+}(T(S))\}$ so, a~basis for the module $\sum U\cdot
T(Q)$ is $\{T(S):\Omega^{+}(T(R))\subsetneq\Omega^{+}(T(S))\}$.
Therefore, $\mathcal{I}(T(R))$ is a~basis for $M(T(R))$.
